\section{A level-set residue estimate for dissipation}
\label{sec:residue}

The local residue controls the full inclusion derivative energy
on sublevel sets of its norm. This gives an actual-flow bound
before constructing a comparison trajectory. The estimate
allows nonreduced defects and coincident line factors.

Fix a coordinate ball $B_R$ about one $p\in\Sigma$, with no
other defect in its closure, and use \eqref{eq:local-koszul}.
The base and line metrics extend smoothly to a neighbourhood
of the closed ball. Set
\begin{equation}\label{eq:jacobian-bound}
 K_p(R)=\sup_{\overline B_R}
       |\dd f\wedge\dd g|_\omega\sqrt{h_M/h_L},
           \qquad m_p=\dim_\CC\OO_{X,p}/(f,g).
\end{equation}
The finite map germ $(f,g)$ has a Jacobian which is not
identically zero, so $K_p(R)>0$ for small $R>0$.
No evolving bundle metric enters this constant.

\begin{lemma}\label{lem:level-set-residue}
For every smooth Hermitian metric $H$ on $V|_{B_R}$ with the
prescribed determinant, let $u=\iota^*H/h_L$ and
$e=|D'_{H\otimes h_L^{-1}}\iota|^2$. Then
\begin{equation}\label{eq:local-residue-bound}
 \int_{B_R}e\,\dd V
       \ge\frac{4\pi^2m_p}{K_p(R)}
                      \left(\min_{\partial B_R}u\right)^2 .
\end{equation}
\end{lemma}

\begin{proof}
Let $\tau$ be the $H$-orthogonal right inverse of $q$ on the
punctured ball, with $\db\tau=\iota\beta$. The form $\beta$
is $\db$-closed, so
$\Omega=\dd f\wedge\dd g\wedge\beta$ is closed.
Its integral over every small outward sphere is
\begin{equation}\label{eq:bm-residue}
                  \int_{S_r}\Omega=4\pi^2m_p .
\end{equation}
Here is the normalization and its independence of the
splitting. The right inverse
$\tau_0=(\bar g,\bar f)^{\mathsf t}/(|f|^2+|g|^2)$
has extension form
\[
 \beta_0=\frac{\bar f\,\dd\bar g-\bar g\,\dd\bar f}
                       {(|f|^2+|g|^2)^2}.
\]
The Bochner--Martinelli degree formula gives
\eqref{eq:bm-residue} for $\beta_0$, with local algebra
length $m_p$; see
\cite[Chapter~II, Section~5, Example~5.5; Chapter~IV,
Section~2, Theorems~2.3--2.4]{Suwa05}.
Any other right inverse differs by $\iota a$, and its
extension form differs by $\db a$. The change in $\Omega$
is the exact form $\dd(a\,\dd f\wedge\dd g)$, whose
sphere integral vanishes.

The orthogonal splitting has induced metrics
$(uh_L,u^{-1}h_M)$. The same Chern calculation as in
\eqref{eq:full-inclusion-energy} gives
\[
                    e=\frac{|\partial u|^2}{u}
                                      +u^3|\beta|_0^2.
\]
The $(1,0)$ part of $\dd u$ wedges to zero with
$\dd f\wedge\dd g$. The product inequality for the remaining
forms, followed by $2AB\le A^2+B^2$, yields
\begin{equation}\label{eq:weighted-residue-form}
 \begin{split}
 |\dd(u^2\Omega)|
 &\le2u|\partial u|\,|\dd f\wedge\dd g|_\omega
                                  \sqrt{h_M/h_L}\,|\beta|_0\\
 &\le K_p(R)
             \left(\frac{|\partial u|^2}{u}+u^3|\beta|_0^2\right).
 \end{split}
\end{equation}
For top-degree complex forms the norm is the modulus of
their coefficient relative to $\dd V$.
In particular $|\dd z_1\wedge\dd z_2|=2$ in the
Euclidean convention; this factor is retained in $K_p(R)$.

Put $b_R=\min_{\partial B_R}u>0$. For a regular value
$0<b<b_R$, the set $U_b=\{x\in B_R:u(x)<b\}$ has compact
closure in $B_R$ and smooth boundary, on which $u=b$.
It contains $p$. Stokes' theorem for $\Omega$ gives
$\int_{\partial U_b}\Omega=4\pi^2m_p$; any component
not containing $p$ contributes zero.

We justify the puncture removal also for $u^2\Omega$.
In fixed smooth coefficient norms put
$|F|=(|f|^2+|g|^2)^{1/2}$. The orthogonal inverse is
\[
                  \tau=\frac{H^{-1}q^*}{qH^{-1}q^*}.
\]
Smooth positivity gives $qH^{-1}q^*\asymp|F|^2$ and
$u\asymp|F|^2$. Differentiating, with constants allowed
to depend on this fixed metric, gives
\[
 |\db\tau|\le C\left(\frac{|\dd F|}{|F|^2}
                                  +\frac1{|F|}\right),
 \qquad
 |\beta|\le C\left(\frac{|\dd F|}{|F|^3}
                                  +\frac1{|F|^2}\right).
\]
Consequently $|u^2\beta|\le C(|F|\,|\dd F|+|F|^2)$
is bounded. Since $\dd f\wedge\dd g$ is bounded, the
integral of $u^2\Omega$ over a radius-$\varepsilon$ sphere
tends to zero. Also \eqref{eq:weighted-residue-form}
makes its exterior derivative integrable, since $e$ is
smooth on the whole ball. Stokes' theorem on
$U_b\setminus B_\varepsilon$ followed by
$\varepsilon\downarrow0$ thus gives
\[
 4\pi^2m_pb^2
      =\left|\int_{U_b}\dd(u^2\Omega)\right|
      \le K_p(R)\int_{U_b}e\,\dd V.
\]
Choose regular values tending upward to $b_R$ to obtain
\eqref{eq:local-residue-bound}. The metric-dependent bounds
above justify puncture removal only; the final constant is
the fixed geometric constant \eqref{eq:jacobian-bound}.
\end{proof}

\begin{theorem}\label{thm:general-upper-rate}
Under the general hypotheses \eqref{eq:extension}, every
solution of \eqref{eq:metric-flow} satisfies
\begin{equation}\label{eq:general-upper-rate}
                           Z(T)=O(T^{-1}).
\end{equation}
If $W$ has a nonreduced point, then
\begin{equation}\label{eq:nonreduced-rate}
       \frac{-Z'(T)}{Z(T)^2}\longrightarrow+\infty,
                          \qquad TZ(T)\longrightarrow0 .
\end{equation}
If every point is reduced, define
\begin{equation}\label{eq:mass-coefficient}
                  C_*=\frac{4\pi^2\sigma}{v_0^2}.
\end{equation}
Then
\begin{equation}\label{eq:reduced-liminf}
 \liminf_{T\to\infty}\frac{-Z'(T)}{Z(T)^2}\ge C_*,
             \qquad \limsup_{T\to\infty}TZ(T)\le C_*^{-1}.
\end{equation}
\end{theorem}

\begin{proof}
Take disjoint fixed coordinate balls about the defects.
The actual metric is smooth at every finite time with
the fixed determinant, so Lemma~\ref{lem:level-set-residue}
applies on each ball. By
Proposition~\ref{prop:qualitative-flow},
$\min_{\partial B_R(p)}u(T)/Z(T)\to1/v_0$.
The exact dissipation identity therefore gives
\begin{equation}\label{eq:fixed-ball-liminf}
 \liminf_{T\to\infty}\frac{-Z'(T)}{Z(T)^2}
       \ge\frac{8\pi^2}{v_0^2}
                          \sum_{p\in\Sigma}\frac{m_p}{K_p(R)}.
\end{equation}
For any such fixed balls the right side is positive.
The derivative of $1/Z$ is thus eventually bounded
below by a positive constant, which proves
\eqref{eq:general-upper-rate}.

At a reduced zero, \eqref{eq:defect-constants} and
the real two-form norm give
\[
 |\dd f\wedge\dd g|_\omega(p)
       =2|\det\partial(f,g)(p)|,\qquad
                 K_p(R)\longrightarrow2\sqrt{d_p}.
\]
Let $R\downarrow0$ after taking the time limit in
\eqref{eq:fixed-ball-liminf}. Since $m_p=1$, this
proves the first assertion in \eqref{eq:reduced-liminf}.
For every $c<C_*$ one has $(1/Z)'\ge c$ for all
sufficiently large times. Integrate and then let
$c\uparrow C_*$ to prove the second assertion.

At a nonreduced zero the derivative of $(f,g)$ cannot
be invertible: the holomorphic inverse theorem would
otherwise give the maximal ideal and length one.
Hence $\dd f\wedge\dd g$ vanishes at that point and
$K_p(R)\to0$. Its summand in
\eqref{eq:fixed-ball-liminf} becomes arbitrarily large.
Given any positive constant, first choose a fixed
radius small enough and then take time sufficiently
large. This proves $(1/Z)'\to+\infty$.
Integration gives $TZ(T)\to0$.
\end{proof}

For reduced points the static metrics
\eqref{eq:balanced-static} satisfy
$2\mathcal E_a/Z_a^2\to C_*$. Thus the coefficient in
\eqref{eq:reduced-liminf} is attained by an explicit
family of smooth compact metrics. This does not yet
establish its attainment along the actual flow. The
remainder of the paper supplies the elliptic correction
and dynamical comparison in the distinct-factor case.
The nonreduced conclusion \eqref{eq:nonreduced-rate}
is an upper-rate statement and does not select a sharp
power or a logarithmic correction.
